\label{ssub:ann}

Artificial Neural Networks (ANNs), and specifically feedforward Multi-Layer Perceptrons (MLPs),
have served as a standard baseline architecture for regression tasks on tabular biomedical data,
trained via backpropagation to minimise a defined loss function such as Mean Squared Error for
continuous targets such as drug sensitivity \cite{Goodfellow2016,Arrieta2020}.
Despite their widespread application, standard MLPs present a well-documented limitation when
applied to high-dimensional genomic data: they are inherently prone to overfitting, particularly
when the number of input features vastly outnumbers the number of biological response samples
\cite{Partin2023}. The GDSC dataset, after preprocessing via the approach established by
Alipour~\cite{Alipour2024}, retains thousands of gene expression features across cancer cell lines
and drug combinations. When an MLP is trained without adequate regularisation, it tends to memorise
training-specific noise rather than capture the true underlying biological relationships between gene
expression profiles and drug response. This manifests as a high generalisation gap, a large
difference between training accuracy and validation accuracy \cite{Ye2024}, and constitutes precisely
the ``Clever Hans'' phenomenon that motivates this research project \cite{Samek2017,Anders2022}.
Pahde et al.~\cite{Pahde2023} demonstrate through their Reveal to Revise (R2R) framework that
state-of-the-art deep neural networks frequently learn spurious correlations embedded in training
data, including medically irrelevant artefacts, and that post-hoc XAI methods applied to such
overfitted models produce explanations that are visually plausible but functionally unreliable.
Although R2R was developed in an image classification context, the underlying mechanism, a model
achieving high training performance by exploiting dataset-specific shortcuts rather than genuine
signal, applies equally to tabular genomic settings, as documented by Partin et
al.~\cite{Partin2023} in the drug response prediction domain. For the purposes of this project, an
intentionally overfitted baseline model serves as a diagnostic tool: its wide generalisation gap
becomes evidence of fragile explanations, which the SAM regularisation intervention in Stage 2 is
designed to correct. This baseline is produced deliberately, by training the network without
regularisation until its generalisation gap stabilises, as detailed in
Section~\ref{sec:methodology}.

\subsubsection{Deep Learning for Cancer Drug Sensitivity Prediction}
\label{ssub:dl-drug-sensitivity}
Beyond the standard MLP baseline, recent work has demonstrated that domain-informed deep learning
architectures significantly outperform generic feedforward networks on drug sensitivity prediction
tasks using GDSC data \cite{Partin2023,Wu2025}. The key limitation of a plain MLP is that it treats
all gene expression features as equally independent, ignoring the rich biological structure known to
govern cancer drug response \cite{Wu2025}, specifically the organisation of genes into signalling
pathways that collectively regulate cancer cell development, proliferation, and resistance mechanisms.
\\
DeepCCDS \cite{Wu2025} is a deep learning framework designed to predict cancer cell drug sensitivity
by characterising cancer driver signals. To overcome the limitations of standard neural networks, it
replaces high-dimensional gene expression data with a prior knowledge network consisting of 38
biologically meaningful signalling pathways \cite{Wu2025}. Alongside these pathway features, the
architecture utilises self-supervised autoencoders to compress gene mutations and drug structures
into low-dimensional embeddings. As demonstrated via t-SNE visualisations (a technique that projects
high-dimensional data onto two dimensions to reveal cluster structure), these refined representations
effectively separate drug-sensitive from drug-insensitive cell populations,
confirming the model captures true biological structure rather than statistical noise \cite{Wu2025}.
DeepCCDS was trained on 319,543 drug–cell line pairs from the GDSC database, demonstrating superior
predictive performance over state-of-the-art methods \cite{Wu2025}. Crucially, its clinical
applicability was validated on solid tumour patient samples from The Cancer Genome Atlas (TCGA) cohort, a large public
repository of real patient tumour genomes, where predicted
LN IC$_{50}$ values successfully distinguished between responsive and non-responsive patient–drug
combinations \cite{Wu2025}. The interpretability properties of DeepCCDS are directly relevant to the
goals of this project. By integrating prior knowledge about cancer driver signals, DeepCCDS exhibits
powerful feature representation capabilities and built-in interpretability. The pathway-level
structure of the architecture makes it possible to trace which biological mechanisms the model
weights most heavily when predicting drug response. This is critically important in the Red XAI
framework: when evolutionary algorithms extract symbolic rules from the model, genes appearing in
those rules can be cross-referenced against the same biological pathways that structured the model's
internal representations, providing a direct route to biological verification via GSEA pathway
enrichment analysis (Section~\ref{ssub:gsea}).

\subsubsection{The Tendency of Deep Models to Overfit Genomic Data and the Clever Hans Problem}
\label{ssub:clever-hans-genomic}
As introduced in Section~\ref{sec:introduction} and contextualised for genomic data in
Section~\ref{ssub:genomics-interpretability}, the Clever Hans problem represents a direct threat to
explanation quality in overfitted models. In the genomic setting specifically, Partin et
al.~\cite{Partin2023} document that the curse of dimensionality, thousands of gene expression
features measured across a comparatively limited number of cell line–drug combinations, creates
conditions under which unregularised models routinely converge on sharp loss landscape minima and
memorise training-specific noise rather than genuine biological signal \cite{Li2018}.
\\
The methodological precedent most relevant to this project is the Reveal to Revise (R2R) framework of
Pahde et al.~\cite{Pahde2023}, which proposes an iterative cycle of bias identification, artefact
localisation, model correction, and re-evaluation. R2R demonstrates that when XAI methods such as LRP
are applied to overfitted models, the resulting explanations attribute high relevance to artefactual
features, producing outputs that are mathematically consistent with the model's behaviour but
medically invalid \cite{Pahde2023,Klauschen2024}. This audit-and-fix philosophy directly anticipates
the Red XAI diagnostic pipeline this project aims to implement through SAM regularisation and
evolutionary rule extraction, with the key distinction that this project applies a regularisation
intervention at the optimisation stage rather than correcting artefacts post hoc. In this project,
the baseline configuration of DeepCCDS \cite{Wu2025}, trained without regularisation to intentionally
exhibit overfitting, serves as the Fragility Baseline against which the SAM-regularised model is
compared in Stage 2. The generalisation gap at this stage is the primary diagnostic indicator of
explanation fragility, as formalised in Section~\ref{sec:problem}.